% AUTHOR: Amlal El Mahrouss
% PURPOSE: WG03: Design of the Nectar Programming Language.

\documentclass[11pt, a4paper]{article}
\usepackage{graphicx}
\usepackage{listings}
\usepackage{amsmath,amssymb,amsthm}
\usepackage{xcolor}
\usepackage{hyperref}
\usepackage[margin=0.5in,top=1in,bottom=1in]{geometry}

\title{Design of the Nectar Programming Language.}
\author{Amlal El Mahrouss\\\texttt{amlal@nekernel.org}}
\date{February 2026}

\definecolor{codegray}{gray}{0.95}
\definecolor{codeblue}{rgb}{0.1,0.1,0.8}
\definecolor{codegreen}{rgb}{0,0.6,0}
\definecolor{codepurple}{rgb}{0.58,0,0.82}

\lstset{
	language=C++,
	backgroundcolor=\color{codegray},
	basicstyle=\footnotesize\ttfamily,
	keywordstyle=\color{codeblue}\bfseries,
	commentstyle=\color{codegreen},
	stringstyle=\color{codepurple},
	numbers=left,
	numberstyle=\tiny\color{gray},
	stepnumber=1,
	numbersep=5pt,
	breaklines=true,
	breakatwhitespace=false,
	frame=single,
	rulecolor=\color{black},
	captionpos=b,
	keepspaces=true,
	showspaces=false,
	showstringspaces=false,
	showtabs=false,
	tabsize=2
}

\begin{document}

\bf \maketitle

\begin{center}
	\rule[0.01cm]{17cm}{0.01cm}
\end{center}

\abstract
{
	Nectar as presented in its primer—is a compiled programming language—designed for systems programming with high-level abstractions.
	It is statically typed—compiled, and supports multi-paradigm programming. This paper covers the mathematical concepts behind the language.
}

\begin{center}
	\rule[1cm]{17cm}{0.01cm}
\end{center}

\section{Definition of the $\lambda$-Execution}

Let a program $\mathbb{P}$ be:

\begin{lstlisting}
extern printf;

const main()
{
    const written := printf("%s:13", "Hello, world!\n");
    return written;
}
\end{lstlisting} 
$\mathbb{P}$ shall execute:

\begin{lstlisting}
$ Hello, world!
\end{lstlisting}
and return variant $written$, now defined as $\mathbb{W}$, upon completion. Such program that we denote as $\Theta$ may be defined as:

\begin{equation}
	\Theta(x) = \lambda x.(P(x))
\end{equation}
where $P(x) = \mathbb{P}$ with an argument of $x$. \\ Let $\Theta(x)$ be defined as the $\lambda$-Execution of a program $\mathbb{P}$.

\section{Definitions}

	\item Nectar: A compiled systems programming language—currently studied in this paper.
	\item $\lambda$-Execution: Formally defined as: $\Theta(x) = \lambda x.(P(x))$.
	\item $\mathbb{W}$: The return variant—an $\lambda$-Execution variable based on the $\lambda$-Execution result.

\section{Conclusion}

We have introduced definitions, and properties which shall be be applied to programming languages, semantics, and other related mathematicial applications.

\section{References}

\begin{enumerate}
  	\item NeKernel.org (2026), \href{https://nekernel.org/nekernel}{https://nekernel.org}
  	\item NeKernel.org Primer (2026), \href{https://nekernel.org/nectar\_primer.html}{https://nekernel.org/nectar\_primer.html}
  	\item El Mahrouss, A. (2026). The Execution Semantics: On Axioms, Domains, and Authority. (v1.0.0). Zenodo. https://doi.org/10.5281/zenodo.18362375
	\item El Mahrouss, A. (2026). The Mathematics of Execution. Zenodo. \\https://doi.org/10.5281/zenodo.18399140
	\item Church, A. (1936). An Unsolvable Problem of Elementary Number Theory. American Journal of Mathematics, 58(2), 345–363. https://doi.org/10.2307/2371045
\end{enumerate}

\end{document}
